% Avaliação complementar: confirmar a posição e a numeração no documento final.
\subsection{Avaliação complementar da decisão integrada}
\label{sec:avaliacao_integracao_sincronizada}

Os quatro classificadores foram reaplicados em horários comuns,
a cada dez minutos, mantendo os modelos treinados e seus tamanhos
de janela. Nesta avaliação complementar, o avanço de inferência
foi de uma leitura. Foram comparados de um a quatro votos e de
uma a três decisões positivas consecutivas. O critério de cada
tarefa foi escolhido somente na validação das cinco rodadas,
priorizando cobertura de pelo menos 90\% dos eventos e, entre
as alternativas elegíveis, maior F1 médio por instante. Quando
nenhuma alternativa atingiu essa cobertura, foi priorizada a
maior cobertura disponível.

A sinalização começou no instante da confirmação, sem preencher
retroativamente a janela. Cada evento foi contado uma única vez
quando houve uma decisão positiva durante seu intervalo. As
primeiras posições, sem janela completa para todos os modelos,
foram registradas como indisponíveis para essa comparação. O
período de teste já havia sido examinado em análises anteriores;
por isso, esta reavaliação tem caráter exploratório. A preparação
com interpolação foi mantida e não permite comprovar operação
causal em tempo real.

\begin{table}[H]
\centering
\caption[Comparação complementar da decisão integrada]
{Comparação complementar da decisão integrada no teste.}
\label{tab:integracao_sincronizada_complementar}
\footnotesize
\setlength{\tabcolsep}{4pt}
\begin{tabular}{llrrrr}
\toprule
 & & \multicolumn{2}{c}{F1 por instante} & \multicolumn{2}{c}{Eventos alcançados} \\
Variável & Categoria & União & Selecionada & União & Selecionada \\
\midrule
PM$_{10}$ & Oscilação & 0.611 & 0.590 & 27/27 & 25/27 \\
PM$_{10}$ & Queda & 0.347 & 0.347 & 26/28 & 26/28 \\
PM$_{2,5}$ & Oscilação & 0.439 & 0.362 & 22/22 & 13/22 \\
PM$_{2,5}$ & Queda & 0.263 & 0.296 & 24/26 & 24/26 \\
Temperatura & Oscilação & 0.449 & 0.455 & 25/25 & 24/25 \\
Temperatura & Queda & 0.269 & 0.229 & 24/26 & 21/26 \\
Umidade & Oscilação & 0.465 & 0.489 & 23/24 & 23/24 \\
Umidade & Queda & 0.277 & 0.280 & 25/27 & 25/27 \\
\bottomrule
\end{tabular}
\fonte{Elaborada pela autora.}
\end{table}

O F1 corresponde à média das cinco rodadas; as contagens de
eventos foram somadas. A união utiliza pelo menos um voto,
sem exigir persistência, nos mesmos horários avaliados pela
regra selecionada. Essas medidas por instante não equivalem
ao F1 por janela das análises individuais.

A concordância entre modelos não melhorou todas as tarefas.
Em Oscilação de PM$_{2,5}$ e Queda de Temperatura, houve perda
de cobertura e de F1. Em Umidade, o F1 aumentou sem redução
da cobertura em relação à união sincronizada, mas Queda ainda
sinalizou cerca de 80\% dos instantes. Assim, reduzir sinalizações
não foi suficiente para demonstrar uma decisão integrada mais
eficaz de forma geral.

% Sugestão para a conclusão: substituir apenas a afirmação correspondente,
% sem duplicar a interpretação acima nem apagar resultados individuais.
% A integração dos classificadores e das regras foi implementada e avaliada.
% A avaliação complementar mostrou que a concordância entre modelos pode
% reduzir sinalizações, mas também deixar eventos sem identificação.
% Portanto, não foi demonstrada melhora uniforme da decisão integrada.
